\ifdefined\gkver\else\def\gkver{student}\fi
\documentclass[\gkver]{gaokaozhenti}
\begin{document}

\chapter{2025年上海}

\begin{enumerate}
\item
%% number: 1
%% paperName: 2025年上海高考第 1 题 12 分
%% typeId: 9
%% type: 综合题
%% chapter: 光学
%% point: 光的干涉
%% method: 无
%% score: 12
%% degree: 824
%% duplicateId: 0
%% body:
\textbf{光学实验}

光学实验是通过可控手段研究光现象及其规律的实践活动，包括光的折射、干涉、衍射、偏振等基本特性验证，以及其他光学性质的研究。
\begin{enumerate}
	\item
	分别用红光和紫光进行杨氏双缝干涉实验时，保持其他条件相同，所得条纹可能是\tkanswer[1cm]{A}。

	\fourchoices[answer=A, ispicture=true, h=1.0cm]
	{figs/01a.png}
	{figs/01b.png}
	{figs/01c.png}
	{figs/01d.png}

	\item
	如图，自然光经过两个环形偏振片，保持 A 不动，偏振片 B 绕圆心所在平面匀速旋转的周期为 $T$，则光屏上相邻两次光强最弱的时间间隔为\tkanswer[1cm]{C}。

	\onepicture[h=2.6cm, align=right]{\includesvg{figs/01e}}

	\fourchoices[answer=C]
	{$2T$}
	{$T$}
	{$0.5T$}
	{$0.25T$}

	\item
	如图，进行“测量玻璃的折射率”实验。在玻璃砖的一侧垂直于纸面插两枚大头针 $a$、$b$，在另一侧透过玻璃砖观察，插入大头针 $c$，使其遮住大头针 $a$、$b$；插入大头针 $d$，使其将\tkanswer[1.5cm]{abc}都挡住，并获得图示光路。测得入射角 $\theta_{1}$ 和折射角 $\theta_{2}$，则玻璃的折射率为\tkanswer[1.5cm]{$\frac{\sin\theta_{1}}{\sin\theta_{2}}$}。

	\onepicture[h=3.0cm, align=right]{\includesvg{figs/01f}}
\end{enumerate}
%% answer: 
\jdanswer{
\begin{enumerate}
	\item
	A

	\item
	C

	\item
	abc，$\frac{\sin\theta_{1}}{\sin\theta_{2}}$
\end{enumerate}
}
%% memo:
\memoanswer{
1．干涉条纹是平行等距明暗相间的条纹，根据 $\Delta x=\frac{L}{d}\lambda$，红光的波长大于紫光，可知红光的条纹间距大于紫光的条纹间距，故选项 A 正确；

2．根据偏振原理，偏振片 B 每转过半周透光强度从最小到最大、再到最小，可知光屏上两个光强最小的时间间隔为 $0.5T$，故选 C。

3．若从 c 侧观察，插入 c 时，应遮住 a、b；插入 d 时，应遮住 c 以及 a、b 的像；

该玻璃的折射率为 $n=\frac{\sin\theta_{1}}{\sin\theta_{2}}$。
}

\item
%% number: 2
%% paperName: 2025年上海高考第 2 题 13 分
%% typeId: 9
%% type: 综合题
%% chapter: 原子和原子核
%% point: 玻尔模型
%% method: 无
%% score: 13
%% degree: 754
%% duplicateId: 0
%% body:
\textbf{量子百年}

2025 年被联合国宣布为国际量子科学与技术年，以致敬量子力学为科学技术领域带来的巨大发展和深远影响。
\begin{enumerate}
	\item
	太阳内部发生的核聚变反应称为 p-p 循环，其核反应方程为 $\ce{4^{1}_{1}H} \to \ce{X} + 2\ce{^{0}_{1}e}$，则 X 是\tkanswer[1cm]{B}。

	\fourchoices[answer=B]
	{氢核}
	{氦核}
	{锂核}
	{铍核}

	\item
	（多选）复色光频率范围为 $5.50\times10^{14} \UHz \sim 6.50\times10^{14} \UHz$，用此光照射下列金属可发生光电效应。由表中各金属的截止频率 $\nu_{0}$ 可知，该金属可能是\tkanswer[1cm]{CDE}。

	\begin{center}
\begin{tabular}{|c|c|c|c|c|c|}
\hline
\multicolumn{6}{|c|}{几种金属的截止频率} \\ \hline
金属 & 锌 & 钙 & 钠 & 钾 & 铷 \\ \hline
$\nu_{0}$（$\times10^{14} \UHz$） & 8.07 & 7.73 & 5.53 & 5.44 & 5.15 \\ \hline
\end{tabular}
\end{center}

	\fivechoices[answer=CDE]
	{锌}
	{钙}
	{钠}
	{钾}
	{铷}

	\item
	氢原子的核外电子以半径 $r$ 绕核做匀速圆周运动，则该电子动量的大小是\tkanswer[3cm]{$e\sqrt{\frac{km}{r}}$}。（电子质量为 $m$，元电荷为 $e$，静电力常量为 $k$）

	\item
	一氢原子在从量子数为 4 的激发态跃迁到量子数为 2 的激发态的过程中，\tkanswer[1cm]{D}。（里德伯常量为 $R$，真空中光速为 $c$）

	\fourchoices[answer=D]
	{吸收一频率为 $Rc\left(\frac{1}{2}-\frac{1}{4}\right)$ 的光子}
	{发射一频率为 $Rc\left(\frac{1}{2}-\frac{1}{4}\right)$ 的光子}
	{吸收一频率为 $Rc\left(\frac{1}{2^{2}}-\frac{1}{4^{2}}\right)$ 的光子}
	{发射一频率为 $Rc\left(\frac{1}{2^{2}}-\frac{1}{4^{2}}\right)$ 的光子}
\end{enumerate}
%% answer: 
\jdanswer{
\begin{enumerate}
	\item
	B

	\item
	CDE

	\item
	$e\sqrt{\frac{km}{r}}$

	\item
	D
\end{enumerate}
}
%% memo:
\memoanswer{
1．设 X 的质量数为 $A$，电荷数为 $Z$。由核反应过程中质量数守恒可得：$4=A+0$；

由电荷数守恒可得：$4=Z+2$；

解得 $A=4$，$Z=2$。对应 $\ce{^{4}_{2}He}$。正确选项为 B。

2．当入射光的频率大于金属的截止频率时，可发生光电效应。与表格数据对比可知可发生光电效应的钠、钾、铷，故选 CDE。

3．电子绕氢原子核做匀速圆周运动，库仑力提供向心力，有

$k\frac{e^{2}}{r^{2}}=m\frac{v^{2}}{r}$

求得 $v=\sqrt{\frac{ke^{2}}{mr}}$

进一步求得电子动量大小 $p=mv=e\sqrt{\frac{km}{r}}$

4．氢原子从高能级向低能级跃迁时会放出光子，光子的能量等于两个能级的能量差。而氢原子能级公式为 $E_{n}=\frac{E_{1}}{n^{2}}$，与量子数 $n$ 的平方成反比，所以从量子数 $n=4$ 跃迁到量子数 $n=2$ 的能量差应有 $\left(\frac{1}{2^{2}}-\frac{1}{4^{2}}\right)$ 的形式。正确选项为 D。

另附严格证明：由广义巴尔末公式 $\frac{1}{\lambda}=R\left(\frac{1}{m^{2}}-\frac{1}{n^{2}}\right)$（$m=1,2,3,\cdots$；$n=m+1,m+2,m+3,\cdots$）、光子能量公式 $E=h\nu=h\frac{c}{\lambda}$，可以推得：$\nu=Rc\left(\frac{1}{m^{2}}-\frac{1}{n^{2}}\right)$。从 $n=4$ 跃迁到 $m=2$ 过程中放出光子的频率为 $Rc\left(\frac{1}{2^{2}}-\frac{1}{4^{2}}\right)$。
}

\item
%% number: 3
%% paperName: 2025年上海高考第 3 题 15 分
%% typeId: 9
%% type: 综合题
%% chapter: 电路
%% point: 电阻定律
%% method: 无
%% score: 15
%% degree: 668
%% duplicateId: 0
%% body:
\textbf{滑动变阻器}

改变滑动变阻器接入电路的方式和阻值可对电路进行调控。小佳用最大阻值为 $10 \UO$ 的滑动变阻器等器材做了一系列实验。
\begin{enumerate}
	\item
	实验室利用如图所示电路进行实验。调节滑动变阻器，当电流表示数 $I_{1}=1 \UA$ 时，电压表示数 $U_{1}=3 \UV$；当电流表示数 $I_{2}=2 \UA$ 时，电压表示数为 $U_{2}=1.5 \UV$，则电源电动势为\tkanswer[1cm]{4.5} $\UV$，内阻为\tkanswer[1cm]{1.5} $\UO$。

	\onepicture[h=2.6cm, align=right]{\includesvg{figs/03a}}

	\item
	小佳用柱形电阻 $X$、电动势为 $3 \UV$ 的电池、滑动变阻器 $R$、电流表、电压表、开关连成电路进行实验，根据实验数据作出通过 $X$ 的电流 $I$ 随 $X$ 两端电压 $U$ 变化的 $I-U$ 图线。实验过程中 $X$ 的长度和横截面积保持不变。

	\onepicture[h=3.5cm, align=right]{\includesvg{figs/03b}}

	\begin{enumerate}
		\item
		若 $X$ 的阻值为 $R_{X}$，$X$ 的材料的电阻率为 $\rho$，根据 $I-U$ 图线可知，当加在 $X$ 两端的电压增大时，\tkanswer[1cm]{B}。

		\fourchoices[answer=B]
		{$R_{X}$ 增大，$\rho$ 增大}
		{$R_{X}$ 减小，$\rho$ 减小}
		{$R_{X}$ 增大，$\rho$ 不变}
		{$R_{X}$ 减小，$\rho$ 不变}

		\item
		当电阻两端电压为 $1.80 \UV$ 时，该电阻的功率约为\tkanswer[1.5cm]{0.225} $\UW$。（保留 3 位有效数字）

		\item
		根据 $I-U$ 图线的测量数据点推测，小佳所采用的实验电路图可能为\tkanswer[1cm]{C}。

		\fourchoices[answer=C, ispicture=true, h=1.6cm]
		{figs/03c.png}
		{figs/03d.png}
		{figs/03e.png}
		{figs/03f.png}
	\end{enumerate}

	\item
	水平放置的平行板电容器连接在图示电路中，电容器极板间距为 $d$，电容为 $C$。移动滑动变阻器的滑片 P，当两极板间电压为 $U$ 时，该电容器所带电荷量为\tkanswer[1.5cm]{$CU$}。此时，一质量为 $m$，电荷量为 $q$ 的带正电粒子以速度 $v_{0}$ 从该电容器左侧两极板间中点位置水平向右射入，恰好从上极板右侧边缘飞出，则两极板长度 $L=$\tkanswer[2.5cm]{$v_{0}d\sqrt{\frac{m}{qU}}$}。（不计带电粒子所受重力和边缘效应）

	\onepicture[h=2.8cm, align=right]{\includesvg{figs/03g}}
\end{enumerate}
%% answer: 
\jdanswer{
\begin{enumerate}
	\item
	4.5，1.5

	\item
	（1）B     （2）0.225     （3）C

	\item
	$CU$，$v_{0}d\sqrt{\frac{m}{qU}}$
\end{enumerate}
}
%% memo:
\memoanswer{
1．根据闭合电路欧姆定律 $E=U+Ir$

代入数据可得 $E=3+r$；$E=1.5+2r$

联立两式解得 $E=4.5 \UV$；$r=1.5 \UO$

2．（1）因为电阻 $I-U$ 图线的斜率为电阻的倒数，随着电阻两端的电压增大，可知斜率在不断增大，故随着电阻两端的电压增大，$R$ 减小，再根据电阻定律 $R=\rho\frac{l}{S}$ 可得 $\rho$ 也在减小，B 选项正确。

（2）从图像上可以看出，当电阻两端电压为 $1.8 \UV$ 时，可读出此时流过电阻的电流为 $0.125 \UA$，根据公式 $P=UI$，代入数据解得 $P=0.225 \UW$。

（3）根据 $I-U$ 图像，实验装置能测得 $U=0$ 附近的数据，因此滑动变阻器应采用分压式接法，而且电流表应测量流过电阻的电流，可推测该实验电路符合条件的只有 C 选项。

3．（1）闭合开关，稳定时，电容器两端的电压为 $U$。根据电容的定义式可知电容器的电荷量为 $Q=CU$。

（2）根据牛顿第二定律得 $qE=q\frac{U}{d}=ma$，得加速度 $a=\frac{qU}{md}$。

粒子恰好从上极板右侧边缘射出，垂直于板的方向有 $\frac{d}{2}=\frac{1}{2}at^{2}$，解得 $t=d\sqrt{\frac{m}{qU}}$

水平方向上，可求得两极板的长度 $L=v_{0}t=v_{0}d\sqrt{\frac{m}{qU}}$
}

\item
%% number: 4
%% paperName: 2025年上海高考第 4 题 13 分
%% typeId: 9
%% type: 综合题
%% chapter: 曲线运动
%% point: 竖直平面圆周运动
%% method: 无
%% score: 13
%% degree: 529
%% duplicateId: 0
%% body:
\textbf{竖直圆轨道}

如图，一光滑圆轨道竖直固定放置，$a$ 为轨道最低点，$b$、$d$ 是与轨道圆心 $O$ 等高的点，$c$ 为轨道最高点，$e$ 为圆弧 $ab$ 的中点。一小滑块以一定的水平初速度从 $a$ 点出发，仅在重力和轨道对其弹力作用下沿圆轨道顺时针运动。

\onepicture[h=3.5cm, align=right]{\includesvg{figs/04a}}
\begin{enumerate}
	\item
	若物块从 $a$ 点运动到 $c$ 点所用时间为 $t_{0}$，则在 $t_{0}/2$ 时刻，物块位于\tkanswer[1cm]{E}。

	\fivechoices[answer=E]
	{$ae$ 之间}
	{$e$ 点}
	{$eb$ 之间}
	{$b$ 点}
	{$bc$ 之间}

	\item
	若质量为 $m$ 的滑块从 $a$ 点以水平速度 $v_{0}$ 出发，恰能运动到 $b$ 点，且用时为 $t$，则此过程中轨道对滑块弹力的冲量大小为\tkanswer[1.5cm]{D}。（重力加速度大小为 $g$）

	\fourchoices[answer=D]
	{$mv_{0}$}
	{$mgt$}
	{$mv_{0}+mgt$}
	{$m\sqrt{v_{0}^{2}+g^{2}t^{2}}$}

	\item
	已知物块质量 $m=0.5 \Ukg$，以初速度 $v_{0}$ 可以沿轨道做完整的圆周运动。图（b）是物块的速率 $v$ 与物块与圆心连线转过的角度 $\theta$ 关系的图像。

	\onepicture[h=3.2cm, align=right]{figs/04b.png}

	\begin{enumerate}
		\item
		求轨道半径 $R$；（计算）

		\item
		当 $\theta=60^{\circ}$ 时，求小球克服重力做功的瞬时功率 $P$。（计算）
	\end{enumerate}
\end{enumerate}
%% answer: 
\jdanswer{
\begin{enumerate}
	\item
	E

	\item
	D

	\item
	（1）0.59 m     （2）23.3 W
\end{enumerate}
}
%% memo:
\memoanswer{
1．物块从 a 点运动到 c 点过程中一直做减速运动，因此物块沿圆弧从 a 点运动到 b 点、从 b 点运动到 c 点，这两个过程虽然路程相同，但前半段路程用时短，因此经过 $t_{0}/2$ 时间物块应运动在 b 点之上，即位于 bc 之间。正确选项为 E。

2．根据动量定理的矢量表达式，有

$I_{\text{合}}=I_{G}+I_{N}=\Delta p$

\onepicture[h=2.5cm, align=right]{figs/04_an.png}

式中 $I_{G}=mgt$，方向竖直向下；$\Delta p=mv_{0}$，方向水平向右，由右图的矢量关系可知：

$I_{N}=\sqrt{I_{G}^{2}+(\Delta p)^{2}}=m\sqrt{g^{2}t^{2}+v_{0}^{2}}$

故选 D。

3．（1）由图像可知，物块在最低点的速度 $v_{0}=6 \Ums$，转过 $60^{\circ}$ 时的速度 $v_{1}=5.5 \Ums$。由动能定理得

$-mg\cdot\frac{R}{2}=\frac{1}{2}mv_{1}^{2}-\frac{1}{2}mv_{0}^{2}$

解得 $R=\frac{v_{0}^{2}-v_{1}^{2}}{g}=\frac{6^{2}-5.5^{2}}{9.8} \Um\approx 0.59 \Um$

（2）物块克服重力做功的瞬时功率

$P=mgv_{1}\sin 60^{\circ}=0.5\times 9.8\times 5.5\times\frac{\sqrt{3}}{2} \UW\approx 23.3 \UW$
}

\item
%% number: 5
%% paperName: 2025年上海高考第 5 题 28 分
%% typeId: 9
%% type: 综合题
%% chapter: 交流电
%% point: LC振荡回路
%% method: 无
%% score: 28
%% degree: 660
%% duplicateId: 0
%% body:
\textbf{特雷门琴}

图示的特雷门琴是一种“不能碰的乐器”，它的两根天线与演奏者的双手构成两个电容器。弹奏者通过在空中移动双手来控制音调和音量。

\onepicture[h=3.2cm, align=right]{\includesvg{figs/05a}}
\begin{enumerate}
	\item
	特雷门琴电路的一部分可看成图示的 $LC$ 回路。将其电磁振荡视为自由振荡。

	\onepicture[h=2.2cm, align=right]{\includesvg{figs/05b}}

	\begin{enumerate}
		\item
		若将手与水平天线构成的电容器近似看成平行板电容器，则当手与水平天线之间的距离减小时，它们构成的电容器的电容\tkanswer[1cm]{A}。

		\threechoices[answer=A]
		{增大}
		{减小}
		{不变}

		\item
		电容器放电完毕瞬间，以下各物理量中达到最大值的有\tkanswer[1cm]{BC}。（多选）

		\fourchoices[answer=BC]
		{电场能}
		{电流}
		{磁场能}
		{电容器两端电压}
	\end{enumerate}

	\item
	特雷门琴的电磁振荡通过扬声器转换为声波。扬声器的音圈绕在永磁体 N 上，其结构的剖面图和正视图分别如图（a）和（b）所示。

	\onepicture[h=3.2cm, align=right]{figs/05c.png}

	\begin{enumerate}
		\item
		A4 音的频率为 $440 \UHz$，$25\UdC$ 时空气中的声速为 $347.6 \Ums$，则 $25\UdC$ 时 A4 音的声波在空气中的波长为\tkanswer[1.5cm]{0.79} $\Um$。

		\item
		当音圈中电流 $I$ 的方向如图（a）所示时，图（a）中音圈所受安培力的方向\tkanswer[1cm]{B}。

		\fourchoices[answer=B]
		{向左}
		{向右}
		{径向向内}
		{径向向外}

		\item
		音圈中一匝线圈的周长为 $2.0 \Ucm$，若音圈所在处的径向磁场磁感应强度大小均为 $0.50 \UT$，音圈中的电流 $i=I_{0}\sin(2\pi ft)$，其中 $I_{0}=0.71 \UA$，$f=200 \UHz$，则该电流的有效值为\tkanswer[1.5cm]{0.50} $\UA$；音圈中一匝线圈受到的安培力的最大值为\tkanswer[2cm]{$7.1\times10^{-3}$} $\UN$。
	\end{enumerate}

	\item
	演奏特雷门琴时，一手控制音调，手与竖直天线间距越小，琴发出的音调越高，另一手控制音量，手与水平天线间距越小，琴发出的音量越小。某特雷门琴输出的一列简谐声波的波形如图中①所示，为使其输出一列波形如②所示的简谐声波，应\tkanswer[1cm]{C}。

	\onepicture[h=2.6cm, align=right]{\includesvg{figs/05d}}

	\fourchoices[answer=C]
	{一手远离竖直天线，另一手靠近水平天线}
	{一手远离竖直天线，另一手远离水平天线}
	{一手靠近竖直天线，另一手靠近水平天线}
	{一手靠近竖直天线，另一手远离水平天线}
\end{enumerate}
%% answer: 
\jdanswer{
\begin{enumerate}
	\item
	（1）A     （2）BC

	\item
	（1）0.79     （2）B     （3）0.50，$7.1\times10^{-3}$

	\item
	C
\end{enumerate}
}
%% memo:
\memoanswer{
1．（1）平行板电容器的电容大小与极板间的距离成反比。当人手靠近天线时，相当于减小了电容器两极板间的距离，从而使电容变大。正确选项为 A。

（2）在 $LC$ 振荡电路中，当电荷量为零时，由 $U=Q/C$ 可知极板的电压为零，电场能也为零；但电流却达到了最大值，而磁场能随电流的增大而增大，故磁场能也达到最大。故选 BC。

2．（1）根据波速、频率和波长的关系，有 $\lambda=\frac{v}{f}=\frac{347.6}{440} \Um=0.79 \Um$。

（2）观察左图，根据左手定则可知线圈所受安培力方向为向右，故选 B。

（3）对于正弦式交变电流，电流的有效值与最大值的关系为 $I=\frac{I_{\mathrm{m}}}{\sqrt{2}}$，则此题中有效值为 $\frac{I_{0}}{\sqrt{2}}=\frac{0.71}{\sqrt{2}} \UA\approx 0.50 \UA$。

将单匝线圈看成由无数段小直导线，则线圈受到的安培力的最大值为 $F=BI_{0}L=0.5\times 0.71\times 0.02 \UN=7.1\times10^{-3} \UN$。

3．由图可知声波由图线①变成图线②频率变大，即音调变高；振幅变小，即音量变轻。而题中告知手靠近竖直天线时，音调变高；靠近水平天线时，音量变轻，可知人手靠近竖直天线、靠近水平天线。故选 C。
}

\item
%% number: 6
%% paperName: 2025年上海高考第 6 题 19 分
%% typeId: 9
%% type: 综合题
%% chapter: 机械振动
%% point: 简谐振动
%% method: 无
%% score: 19
%% degree: 410
%% duplicateId: 0
%% body:
\textbf{再生制动与气囊减震}

再生制动与气囊减震相辅相成，共同提升车辆性能。再生制动通过能量回收提升能源利用效率，制动系统优先使用再生制动，不足部分由机械制动补充。气囊减震提升了乘坐舒适性和安全性。

\onepicture[h=3.2cm, align=right]{figs/06a.png}
\begin{enumerate}
	\item
	一电动汽车采用笼式感应电机，再生制动时电机转为发电机模式。如图（a）所示的笼式电机在磁感应强度为 $B$ 的匀强磁场中绕固定轴 $OO'$ 转动，磁场方向与各导体条垂直。当电机转到如图（b）所示位置时，感应电动势最大的导体条为\tkanswer[1.5cm]{A}。

	\fourchoices[answer=A]
	{$a_{1}$ 和 $a_{2}$}
	{$b_{1}$ 和 $b_{2}$}
	{$c_{1}$ 和 $c_{2}$}
	{$d_{1}$ 和 $d_{2}$}

	\item
	一质量 $m=1500 \Ukg$ 的电动汽车沿坡度 $\theta=4.8^{\circ}$ 的斜坡匀减速下坡。从 $v_{0}=72 \Ukmh$ 减速至 $v_{t}=18 \Ukmh$ 用时 $t=5 \Us$。若这 $5 \Us$ 内系统完全使用再生制动，且该车在行驶过程中受到的除制动力以外的其余阻力的大小 $f=500 \UN$，求这 $5 \Us$ 内该车的位移大小 $x$ 和再生制动提供的制动力大小 $F$。（计算）

	\item
	气囊减震技术中的气囊可简化为图示水平固定的气缸，质量为 $M$、横截面积为 $S$ 的光滑活塞将一定质量的理想气体封闭在气缸内。开始时气缸内气体压强为 $p_{0}$、体积为 $V_{0}$，活塞静止于平衡位置 $O$。外界大气压强恒为 $p_{0}$。以 $O$ 为原点、水平向右为正方向建立 $x$ 轴。

	\onepicture[h=2.6cm, align=right]{\includesvg{figs/06b}}

	\begin{enumerate}
		\item
		若缸内气体经历等温过程，求活塞在平衡位置 $O$ 右侧、位移为 $x$ 时内外气体对活塞的压力的合力大小 $F$。（计算）

		\item
		若缸内气体经历等温过程，活塞由平衡位置 $O$ 沿 $x$ 轴正方向移动微小位移并由静止释放，证明活塞近似做简谐振动（论证）。与振子质量为 $m$、弹簧劲度系数为 $k$ 的弹簧振子频率公式 $f_{\text{弹}}=\frac{1}{2\pi}\sqrt{\frac{k}{m}}$ 类比，可知该活塞做简谐振动的频率 $f_{1}=$\tkanswer[2.5cm]{$\frac{S}{2\pi}\sqrt{\frac{p_{0}}{MV_{0}}}$}。【提示：当 $|y|\ll 1$ 时，$\frac{y}{y+1}\approx y$】

		\item
		若气缸和活塞导热性均不佳，活塞移动微小位移并由静止释放，在短时间内活塞的运动仍可视为简谐运动，其频率为 $f_{2}$。

		\begin{steps}
			\item
			与等温情况下活塞近似做简谐运动的频率 $f_{1}$ 相比，\tkanswer[1cm]{C}。

			\threechoices[answer=C]
			{$f_{1}>f_{2}$}
			{$f_{1}=f_{2}$}
			{$f_{1}<f_{2}$}

			\item
			对上述①中的选择作出解释。（简答）
		\end{steps}
	\end{enumerate}
\end{enumerate}
%% answer: 
\jdanswer{
\begin{enumerate}
	\item
	A

	\item
	（1）$x=62.5 \Um$     （2）$F=5230 \UN$

	\item
	（1）$F=\frac{p_{0}S^{2}}{V_{0}+xS}x$     （2）$f=\frac{S}{2\pi}\sqrt{\frac{p_{0}}{MV_{0}}}$     （3）C。根据热力学第一定律，若气缸绝热，则当气体体积增大时，气体对外做功，内能减小，气体温度降低，则压强 $p$ 变得更小，则 $F_{\text{回}}=-(p_{0}-p)S$ 变得更大，$k$ 也随之变大，由 $f=\frac{1}{2\pi}\sqrt{\frac{k}{M}}$ 可知 $f$ 也变大，即 $f_{1}<f_{2}$。
\end{enumerate}
}
%% memo:
\memoanswer{
1．由图可知，磁场方向竖直向下，图中各点的线速度都沿切线方向，而此时 $a_{1}$ 和 $a_{2}$ 点速度方向与磁场方向垂直，产生的感应电动势最大。故选 A。

2．解：$v_{0}=72 \Ukmh=20 \Ums$，$v_{t}=18 \Ukmh=5 \Ums$

车沿斜坡匀减速运动     $x=\frac{v_{0}+v_{t}}{2}t=\frac{20+5}{2}\times 5 \Um=62.5 \Um$

车沿斜坡匀减速运动，设其加速度大小为 $a$

$v_{t}-v_{0}=at$

据牛顿第二定律可知 $F+f-mg\sin\theta=ma$

\onepicture[h=2.6cm, align=right]{figs/06_an.png}

联立解得：$F=m\frac{v_{t}-v_{0}}{t}-f+mg\sin\theta$

$=(1500\times\frac{20-5}{5}-500+1500\times 9.8\times\sin 4.8^{\circ}) \UN=5230 \UN$

3．（1）设活塞偏离平衡位置的位移为 $x$ 时，气缸内气体压强为 $p$。等温情况下气缸内气体满足玻意耳定律

$p_{0}V_{0}=p(V_{0}+xS)$

气缸内外压力差 $F=(p_{0}-p)S$

解得 $F=\frac{p_{0}S^{2}}{V_{0}+xS}x$

（2）设 $x$ 方向为正方向，则此时活塞所受合力作为回复力

$F_{\text{回}}=-F=-\frac{p_{0}S^{2}}{V_{0}+xS}x=-\frac{\frac{S}{V_{0}}x}{1+\frac{S}{V_{0}}x}p_{0}S$

由题意：若 $|y|\ll 1$，$\frac{y}{y+1}\approx y$；当 $x$ 很微小时，$\frac{S}{V_{0}}x\ll 1$，所以 $\frac{\frac{S}{V_{0}}x}{1+\frac{S}{V_{0}}x}\approx\frac{S}{V_{0}}x$；即

$F_{\text{回}}\approx -p_{0}S\frac{S}{V_{0}}x=-kx$

即活塞的振动可视为简谐振动，其中比例系数 $k=\frac{p_{0}S^{2}}{V_{0}}$

频率 $f=\frac{1}{2\pi}\sqrt{\frac{k}{M}}=\frac{S}{2\pi}\sqrt{\frac{p_{0}}{MV_{0}}}$

（3）根据热力学第一定律，若气缸绝热，则当气体体积增大时，气体对外做功，内能减小，气体温度降低，则压强 $p$ 变得更小，则 $F_{\text{回}}=-(p_{0}-p)S$ 变得更大，$k$ 也随之变大，由 $f=\frac{1}{2\pi}\sqrt{\frac{k}{M}}$ 可知 $f$ 也变大，即 $f_{1}<f_{2}$。故选 C。
}

\end{enumerate}

\end{document}
